\documentclass[10pt,a4paper]{amsart}
\usepackage{amsmath,amssymb,amsthm,graphicx}
% For commutative diagrams we use Xy-pic
\usepackage[all]{xy}

% We use 2cell for 2-commutative diagrams.
\xyoption{2cell}
% Theorem environments.
%
\theoremstyle{plain}
\newtheorem{theorem}[subsection]{Theorem}
\newtheorem{proposition}[subsection]{Proposition}
\newtheorem{lemma}[subsection]{Lemma}

\theoremstyle{definition}
\newtheorem{definition}[subsection]{Definition}
\newtheorem{example}[subsection]{Example}
\newtheorem{exercise}[subsection]{Exercise}
\newtheorem{situation}[subsection]{Situation}

\theoremstyle{remark}
\newtheorem{remark}[subsection]{Remark}
\newtheorem{remarks}[subsection]{Remarks}

\numberwithin{equation}{subsection}

% Macros
%
\def\lim{\mathop{\mathrm{lim}}\nolimits}
\def\colim{\mathop{\mathrm{colim}}\nolimits}
\def\Spec{\mathop{\mathrm{Spec}}}
\def\Hom{\mathop{\mathrm{Hom}}\nolimits}
\def\Ext{\mathop{\mathrm{Ext}}\nolimits}
\def\SheafHom{\mathop{\mathcal{H}\!\mathit{om}}\nolimits}
\def\SheafExt{\mathop{\mathcal{E}\!\mathit{xt}}\nolimits}
\def\Sch{\mathit{Sch}}
\def\Mor{\mathop{\mathrm{Mor}}\nolimits}
\def\Ob{\mathop{\mathrm{Ob}}\nolimits}
\def\Sh{\mathop{\mathit{Sh}}\nolimits}
\def\NL{\mathop{N\!L}\nolimits}
\def\CH{\mathop{\mathrm{CH}}\nolimits}
\def\proetale{{pro\text{-}\acute{e}tale}}
\def\etale{{\acute{e}tale}}
\def\QCoh{\mathit{QCoh}}
\def\Ker{\mathop{\mathrm{Ker}}}
\def\Im{\mathop{\mathrm{Im}}}
\def\Coker{\mathop{\mathrm{Coker}}}
\def\Coim{\mathop{\mathrm{Coim}}}

% Boxtimes
%
\DeclareMathSymbol{\boxtimes}{\mathbin}{AMSa}{"02}

%
% Macros for moduli stacks/spaces
%
\def\QCohstack{\mathcal{QC}\!\mathit{oh}}
\def\Cohstack{\mathcal{C}\!\mathit{oh}}
\def\Spacesstack{\mathcal{S}\!\mathit{paces}}
\def\Quotfunctor{\mathrm{Quot}}
\def\Hilbfunctor{\mathrm{Hilb}}
\def\Curvesstack{\mathcal{C}\!\mathit{urves}}
\def\Polarizedstack{\mathcal{P}\!\mathit{olarized}}
\def\Complexesstack{\mathcal{C}\!\mathit{omplexes}}
% \Pic is the operator that assigns to X its picard group, usage \Pic(X)
% \Picardstack_{X/B} denotes the Picard stack of X over B
% \Picardfunctor_{X/B} denotes the Picard functor of X over B
\def\Pic{\mathop{\mathrm{Pic}}\nolimits}
\def\Picardstack{\mathcal{P}\!\mathit{ic}}
\def\Picardfunctor{\mathrm{Pic}}
\def\Deformationcategory{\mathcal{D}\!\mathit{ef}}
\usepackage{cleveref}
\crefname{originallemma}{original lemma}{original lemmas}
\Crefname{originallemma}{Original lemma}{Original lemmas}
\makeatletter
\newlabel{more-morphisms-lemma-stratify-flat-fp-qf}{{07RZ}{source}{}{sourcepointer-1}{}}
\newlabel{more-morphisms-lemma-stratify-flat-fp-qf@cref}{{[originallemma][0][]07RZ}{[1][0][]source}{}{}{}}
\makeatother
\crefname{originallemma}{original lemma}{original lemmas}
\Crefname{originallemma}{Original lemma}{Original lemmas}
\makeatletter
\newlabel{morphisms-lemma-image-proper-scheme-closed}{{01W6}{source}{}{sourcepointer-2}{}}
\newlabel{morphisms-lemma-image-proper-scheme-closed@cref}{{[originallemma][0][]01W6}{[1][0][]source}{}{}{}}
\makeatother
\crefname{originallemma}{original lemma}{original lemmas}
\Crefname{originallemma}{Original lemma}{Original lemmas}
\makeatletter
\newlabel{schemes-lemma-characterize-separated}{{01KP}{source}{}{sourcepointer-3}{}}
\newlabel{schemes-lemma-characterize-separated@cref}{{[originallemma][0][]01KP}{[1][0][]source}{}{}{}}
\makeatother
\title{Stratifying the base so that an etale morphism becomes finite etale}
\author{The Stacks Project Authors}
\date{stacks-project:etale-cohomology.tex:local-snapshot-20261011; GNU FDL 1.2 or later}
\begin{document}
\maketitle
\noindent\textit{Editorial note.} All delivered mathematical text is the author's own, copied from the pinned Stacks Project chapter file etale-cohomology.tex; no statement or proof step was rewritten, reordered or removed.\par\medskip
\noindent\textit{Editorial note.} The proof is the author's complete argument for both cases and for the combination of the two stratifications.\par\medskip
\noindent\textit{Editorial note.} Every labelled reference occurring inside the counted spans is carried in assembly.external\_references as a non-load-bearing navigation pointer (same-chapter ones at their definition site, cross-chapter ones at the invocation inside the counted span), so no load-bearing mathematics is replaced by a stub. No bibliography citation occurs inside the counted spans.\par\medskip
\noindent\textit{Editorial note.} The checked chapter list for same-name collisions is the set of stacks chapters staged in this repository: chow, cohomology, divisors, etale-cohomology, flat, more-algebra, more-morphisms, morphisms, spaces-more-morphisms, varieties together with this card's chapter etale-cohomology.\par\medskip
\noindent\textit{Editorial note.} The preamble is the chapter's own native preamble excerpt already used for this source: the Xy-pic setup (native preamble lines 14-18) and the native theorem declarations with the mathematical macros (native preamble lines 154-217).\par\medskip
\noindent\textit{Allowed standard prerequisites (editorial).}\begin{itemize}
\item Constructible locally closed subschemes and partitions of a quasi-compact quasi-separated scheme. Same chapter, the notion of a constructible locally closed subscheme and the existence of finite partitions into such subschemes. Supplies the strata of the partition produced in the proof.
\item Finiteness of an etale morphism over a suitable locally closed stratum. Same chapter together with the chapter Morphisms, the description of an etale morphism as finite etale after restricting to a suitable locally closed subscheme, together with the separated and quasi-finite reductions. Supplies the construction of the partition stratum by stratum.
\end{itemize}\medskip
\ifdefined\hypertarget\hypertarget{sourcepointer-1}{}\fi
\noindent\textit{Source reference.} Original lemma 07RZ: sources/stacks-etale-cohomology/pinned-original-source.tex, source lines 12044--12044. Navigation pointer invoked at proof line 12044; the counted span contains the pointer only and the cited result is not reproduced inside it, so the pointer is registered as a non-load-bearing navigation stub in the established form for Stacks material. No bibliography citation is involved.\par\medskip
\ifdefined\hypertarget\hypertarget{sourcepointer-2}{}\fi
\noindent\textit{Source reference.} Original lemma 01W6: sources/stacks-etale-cohomology/pinned-original-source.tex, source lines 12055--12055. Navigation pointer invoked at proof line 12055; the counted span contains the pointer only and the cited result is not reproduced inside it, so the pointer is registered as a non-load-bearing navigation stub in the established form for Stacks material. No bibliography citation is involved.\par\medskip
\ifdefined\hypertarget\hypertarget{sourcepointer-3}{}\fi
\noindent\textit{Source reference.} Original lemma 01KP: sources/stacks-etale-cohomology/pinned-original-source.tex, source lines 12058--12058. Navigation pointer invoked at proof line 12058; the counted span contains the pointer only and the cited result is not reproduced inside it, so the pointer is registered as a non-load-bearing navigation stub in the established form for Stacks material. No bibliography citation is involved.\par\medskip
\begin{lemma}
\label{lemma-etale-stratified-finite}
Let $U \to X$ be an \'etale morphism of quasi-compact and quasi-separated
schemes (for example an \'etale morphism of Noetherian schemes). Then there
exists a partition $X = \coprod_i X_i$ by constructible locally closed
subschemes such that $X_i \times_X U \to X_i$ is finite \'etale for all $i$.
\end{lemma}

\begin{proof}
If $U \to X$ is separated, then this is
More on Morphisms, Lemma \ref{more-morphisms-lemma-stratify-flat-fp-qf}.
In general, we may assume $X$ is affine. Choose a finite affine open
covering $U = \bigcup U_j$. Apply the previous case to all the morphisms
$U_j \to X$ and $U_j \cap U_{j'} \to X$ and choose a common
refinement $X = \coprod X_i$ of the resulting partitions.
After refining the partition further we may assume $X_i$ affine as well.
Fix $i$ and set $V = U \times_X X_i$. The morphisms
$V_j = U_j \times_X X_i \to X_i$ and
$V_{jj'} = (U_j \cap U_{j'}) \times_X X_i \to X_i$ are finite \'etale.
Hence $V_j$ and $V_{jj'}$ are affine schemes and $V_{jj'} \subset V_j$
is closed as well as open (since $V_{jj'} \to X_i$ is proper, so
Morphisms, Lemma \ref{morphisms-lemma-image-proper-scheme-closed}
applies). Then $V = \bigcup V_j$ is separated because
$\mathcal{O}(V_j) \to \mathcal{O}(V_{jj'})$ is surjective, see
Schemes, Lemma \ref{schemes-lemma-characterize-separated}.
Thus the previous case applies to $V \to X_i$ and we can further
refine the partition if needed (it actually isn't but we don't
need this).
\end{proof}

\end{document}
